\paragraph{\trLinearization。}\label{sub:linearization}
\eqref{eq:inhomogeneous-equation} 的解集 $\scr{M}_S(\{y_\zeta\})$ 可纳入一个标准的
函数空间框架。取定 $2<p<\infty$。令 $\scr{B}_S$ 为由映射 $u\colon S\to M$ 组成的%
\trBanachManifold，这些映射是局部 $W^{1,p}$ 的，并且在诸\trStriplikeEnd{}上以
$W^{1,p}$ 的意义\trConvergence{}到某个族 $\{y_\zeta\}$。这个\trManifold{}在一点
$u$ 处的\trTangentSpace{}恰好是 $W^{1,p}(S,E,F)$，其中 $E=u^*TM$，而对每个%
\trConnectedComponent{}$C\subset\partial S$ 有 $F|C=u^*TL_C$。给 $E$ 配备%
\trAlmostComplex{}结构 $J_E(z)=J(z,u(z))$，这里 $J$ 是我们的\trPerturbationDatum{}%
的一部分。有一个自然的\trBanach{}\trVectorBundle{}$\scr{E}_S\to\scr{B}_S$，其诸%
\trFiber{}是空间 $L^p(S,\Omega^{0,1}_S\otimes u^*TM)$。可以把 $\bar\partial-\nu$
想成这个丛的一个\trSmooth\trSection，而它的零点集显然就是诸\trModuliSpace{}%
$\scr{M}_S(\{y_\zeta\})$ 的并（从而赋予它们一个适当的\trTopology）。\trModuliSpace{}%
中一点处的\trLinearization，记作
\begin{equation}\label{eq:linearization}
D_{S,u}\colon (T\scr{B}_S)_u\longrightarrow(\scr{E}_S)_u,
\end{equation}
是一个\trNondegenerate\trCauchyRiemann{}算子，因而是\trFredholm{}的。我们称 $u$ 是%
\textbf{\trRegular}（regular）的，若 $D_{S,u}$ 为\trSurjective。若这对一切
$u\in\scr{M}_S(\{y_\zeta\})$ 都成立，我们就称该\trModuliSpace{}本身是\trRegular{}的。
这蕴含它是\trSmooth{}的且具有相应的维数，后者在每一点处由 $D_{S,u}$ 的%
\trFredholm{}\trIndex{}给出。最后，若这对一切 $\{y_\zeta\}$ 都成立，我们便说%
\trPerturbationDatum{}$(K,J)$ 本身是\trRegular{}的（同样的术语也用于\trFloerDatum）。
要把 $D_{S,u}$ 与上文所考虑的那类显式算子 $\bar\partial_\nabla$ 等同起来，需给 $E$
配备从 $\phi$ \trPullback{}来的\trSymplecticStructure。这个\trConnection{}$\nabla$
并不唯一，因为进入算子定义的只是它的 $(0,1)$-分量；不过，一种可能的选取方式如下。
设 $\tilde u\colon S\to S\times M$ 是 $u$ 的图。取定一个
$Z\in C^\infty(S\times M,TS^*\otimes TM)$，使得 $Z(z,u(z))=du(z)$，并且在每一处
$(Z-Y)^{0,1}=0$。对任何一个这样的 $Z$，都能配以 $E$ 上的一个\trConnection{}$\nabla$。
为此，取 $\widetilde X\in C^\infty(S\times M,TM)$，记其经 $\tilde u$ 的\trPullback{}%
为 $X\in C^\infty(S,E)$，并令
\begin{equation}\label{eq:connection-definition}
\nabla X=\tilde u^*[Z,\widetilde X].
\end{equation}
为看出这与 $\widetilde X$ 的选取无关、而且的确定义了一个\trConnection，任取
$TM\to S\times M$ 上一个无\trTorsion{}的\trConnection{}$\widetilde\nabla^0$，并记其经
$\tilde u$ 的\trPullback{}为 $\nabla^0$。借助它可以把\eqref{eq:connection-definition}
改写为
\begin{equation*}
\nabla_\xi X=\nabla^0_\xi X-\tilde u^*\bigl(\widetilde\nabla^0_X\langle Z,\xi\rangle\bigr)
\end{equation*}
对一切 $\xi\in C^\infty(S,TS)$（其中 $\langle Z,\xi\rangle\in
C^\infty(S\times M,TM)$ 是自然\trPairing{}$TS\otimes(T^*S\otimes TM)\to TM$）。
这表明 $\nabla$ 具有 $\nabla^0$ 加上一个零阶项的形状，这个零阶项依赖于 $Z$ 沿 $u$
的图的导数。接下来，利用 $S$ 上的局部\trHolomorphic{}坐标 $z=s+it$ 把事情写得更
显式，可以算出 $\nabla$ 的 $(0,1)$ 部分为
\begin{equation*}
\begin{split}
(\nabla X)^{0,1}
&=\tfrac12\bigl(\nabla^0_{\partial_s}X+J\nabla^0_{\partial_t}X
-\widetilde\nabla^0_X\langle Z,\partial_s\rangle
-J\widetilde\nabla^0_X\langle Z,\partial_t\rangle\bigr)\\
&=\tfrac12\bigl(\nabla^0_{\partial_s}X+J\nabla^0_{\partial_t}X
+\tilde u^*(\widetilde\nabla^0_XJ)\,du(\partial_t)\bigr)\\
&\qquad-\tfrac12\,\tilde u^*\nabla^0_X\bigl(Y(\partial_s)+JY(\partial_t)\bigr).
\end{split}
\end{equation*}
这的确是\trPseudoholomorphicMap{}理论中\trLinearization{}算子的标准形式，只是添上了
最后一个被加项，以便把\eqref{eq:inhomogeneous-equation} 中的\trInhomogeneous{}项
考虑进来；于是 $D_{S,u}=\bar\partial_\nabla$。为了满足\trAdmissible{}性条件，我们
这样选取 $Z$：如果 $s+it$ 是某个\trStriplikeEnd{}$\epsilon_\zeta\colon Z^\pm\to S$
上的标准坐标，则 $Z(\partial_s)$ 随 $s\to\pm\infty$ 指数式地趋于 $0$（在 $C^2$
\trTopology{}下）。在这种情形，$(E,\phi_E,J_E,F)$ 与 $\nabla$ 的\trAsymptotic{}结构
由如下\trLimitingDatum{}描述：\trVectorBundle{}$E_\zeta=y_\zeta^*TM$，带有从
$\phi$ 继承的\trSymplecticStructure，以及\trAlmostComplex{}结构
$J_{E_\zeta}(t)=J_\zeta(t,y(t))$\footnote{译注：原书此处印作
$J_\zeta(t,y(t))$，且下文 $\Lambda_{\zeta,k}=T(L_{\zeta,k})_{y(k)}$ 中的 $y(k)$
亦不带 $\zeta$ 下标；疑均应为 $y_\zeta(t)$、$y_\zeta(k)$（承前
$E_\zeta=y_\zeta^*TM$ 与“沿其\trIntegralCurve{}$y_\zeta$”之语境）；此处照原书
排印，存疑备案。——译者注}；\trLagrange{}子空间
$\Lambda_{\zeta,k}=T(L_{\zeta,k})_{y(k)}$；以及由 $H_\zeta$ 的“流”沿其%
\trIntegralCurve{}$y_\zeta$ \trLinearization{}而得的\trSymplectic\trConnection{}%
$\nabla_\zeta$。
\begin{remark}\label{thm:linearization-special-cases}
在一些特殊情形下，\trLinearization{}算子有一个更直接的描述。设 $J$ 与 $z\in S$ 无关
且可积，并设 $Y=0$。此时 $TM$ 及其\trPullback{}$E$ 自然地是\trHolomorphic{}丛，而
$D_{S,u}$ 就是与之相配的\trCanonical{}\trDolbeault{}算子。另一种这样的情形是：
$u\equiv x$ 是常值的，且 $Y$ 连同其一阶导数都在点 $x$ 处为零。此时二阶导数 $D^2Y$
可视为 $\Omega^1(S,\mathrm{sym}^2(TM_x)^\vee)$ 中的一个元素，从而可视为在平凡丛
$E=S\times TM_x$ 上给出了一个\trSymplectic\trConnection。于是 $D_{S,u}$ 就是这个%
\trConnection{}的 $(0,1)$ 部分，其中所取的\trComplexStructure{}为 $J_E(z)=J(z,x)$。
\end{remark}
